%% title: Tier-based strict locality
%% category: misc

Reducing long-distance agreement to local interactions is essentially the project of \textbf{tier-based strictly local (TSL) syntax,} a framework arguing that most if not all syntactic phenomena belong to a specific sub-regular computational class (‘tier-based strictly local’). The idea behind TSL is that one divides syntactic phenomena into ‘tiers’ for each feature type, like a person tier and a case tier, and then for each tier, removes nodes of a syntax tree that are not relevant to that tier. Finally, one enforces local (rather than long-distance) constraints on these reduced trees.

In this squib I argue that TSL syntax gets a lot of things right, and also some things wrong. It characterizes observed syntactic phenomena more narrowly than any other computational class, yet it still undergenerates in some places, and overgenerates in others. In some cases, overgeneration must be punted to other, independently motivated parts of the grammar (the lexicon, or selectional restrictions), but in others, both overgeneration and undergeneration can be fixed by adding \textbf{derivationality} to the TSL theory. In other words, we need to view TSL not as a theory which accepts or rejects completed syntax trees, but instead as a theory that proposes specific syntactic operations which obligatorily happen at specific times.
